% ============================================================================
% Chapter 1: Introduction
% ============================================================================
\chapter{Introduction}
\label{ch:introduction}

\section{Background and Motivation}

Advanced Persistent Threats (APTs) are stealthy, multi-stage cyber campaigns that compromise
enterprise systems over extended periods. Unlike commodity malware that executes a single payload
and terminates, APTs establish persistent footholds, conduct lateral movement, escalate
privileges, and exfiltrate sensitive data over weeks or months. Because APTs execute
benign-looking commands and proceed through multi-hop lateral movements, traditional
signature-based security systems are ineffective against them.

Host audit logging and provenance graph analysis have emerged as powerful paradigms for APT
detection. A provenance graph models system events (e.g., file reads, process spawns, network
flows) as directed edges between system entities (processes, files, network sockets), preserving
the complete causal history of system execution.

Despite their potential, existing state-of-the-art provenance-based IDS---such as
OCR-APT~\cite{ocrapt2025}, TraceCluster~\cite{tracecluster2026}, and
StageFinder~\cite{stagefinder2026}---share a \textbf{fatal limitation: the clean-data
assumption}. They require a pristine benign period or fully labeled training dataset. This
assumption fails in practice for four fundamental reasons:

\begin{enumerate}[leftmargin=2cm]
\item \textbf{Stealthy Contamination}: Attackers may already be active when auditing begins,
    corrupting the normal training data.
\item \textbf{Concept Drift}: Enterprise systems naturally evolve due to software updates,
    new administrative activities, and infrastructure changes, making static models obsolete.
\item \textbf{Explosive Log Volumes}: The DARPA OpTC dataset alone contains over 17 million
    provenance edges. Manual labeling is economically intractable.
\item \textbf{False Positive Fatigue}: Static anomaly scoring lacks statistical controls on
    alarm rates, overwhelming analysts with false alerts.
\end{enumerate}

These challenges collectively form the \textit{clean-data deployment bottleneck}---a
fundamental barrier that prevents provenance-based IDS from transitioning from laboratory
research to practical enterprise deployment.

\section{Research Problem}

The central research question addressed in this thesis is:

\begin{quote}
\textit{Can we detect APT attacks by identifying violations of causal invariances discovered
directly from raw, unlabeled, and potentially contaminated system logs---without requiring any
benign training data and with provable false-alarm controls?}
\end{quote}

This question encapsulates three interlinked sub-problems:
\begin{enumerate}
\item \textbf{RQ1}: How can stable causal mechanisms be discovered online from streaming,
    unlabeled provenance data without assuming the training data is clean?
\item \textbf{RQ2}: How can deviations from these mechanisms be projected into a rich
    behavioral state representation for multi-perspective threat reasoning?
\item \textbf{RQ3}: How can detection thresholds be calibrated online to provide empirical
    false-alarm control under concept drift and training contamination?
\end{enumerate}

\section{Proposed Solution: CBSR}

To address these challenges, we propose \textbf{Causal Behavioral State Reasoning (CBSR)},
a framework that detects APT attacks using uncurated streaming data with no manual labels and
no offline retraining.

Our core insight is that \textbf{causal relationships are more stable than statistical
correlations}. While statistical distributions of events vary over time due to concept drift,
the underlying causal equations governing system execution remain invariant. An APT attack
violates these learned causal mechanism equations by injecting novel relationships (e.g., a
Word document spawning PowerShell) or altering the dependencies between processes and files.

Rather than relying solely on noisy individual causal residuals, CBSR projects SCM outputs
into a formal 12-dimensional \emph{Behavioral State Space} and analyzes this state space
from three complementary perspectives: instantaneous mechanism violations ($R_L$), temporal
sequence patterns ($R_T$), and behavioral manifold geometry ($R_B$). Their outputs are fused
via a calibrated meta-learner and constrained by online conformal prediction.

\section{Thesis Contributions}

The key contributions of this thesis are:

\begin{enumerate}
\item \textbf{Causal Behavioral State Space}: A formal 12-dimensional state representation
    derived from online SCM invariants, organized into five coordinate groups (Structural
    Consistency, Activity Dynamics, Behavioral Complexity, Mechanism Violation, Intervention
    Locality) to capture diverse manifestations of APT behavior.

\item \textbf{Causal Intervention Score (CIS)}: A novel graph-boundary metric grounded in
    spectral graph theory that quantifies multi-hop violation propagation, providing principled
    separation between localized administrative noise and multi-stage APT lateral movement.

\item \textbf{Multi-Perspective Parallel Reasoning}: A three-operator architecture ($R_L$,
    $R_T$, $R_B$) that analyzes the same behavioral state from local, temporal, and
    behavioral perspectives simultaneously, resolving the model-mismatch failures of pure
    causal detectors.

\item \textbf{Calibrated Risk Inference with Conformal Guarantees}: A stacked meta-fusion
    layer with online conformal prediction that delivers distribution-free FPR bounds and
    adapts to concept drift via change-point-triggered queue flushing.

\item \textbf{Multi-Benchmark Validation}: Evaluation on five diverse benchmarks (DARPA OpTC,
    BETH, StreamSpot, TC3, NODLINK) demonstrating AUC of 0.9628 on real enterprise Windows
    logs, AUC 0.9901 on Linux sysdig telemetry, AUC 0.9837/0.9784 on simulated APT datasets,
    and a 13$\times$ runtime speedup enabling real-time streaming deployment.
\end{enumerate}

\section{Thesis Organization}

The remainder of this thesis is organized as follows:
\begin{itemize}
\item \textbf{Chapter~\ref{ch:literature}}: Reviews provenance-based IDS, causal anomaly
    detection methods, and conformal prediction, identifying the research gaps that CBSR
    addresses.
\item \textbf{Chapter~\ref{ch:threat_model}}: Defines the adversary model and security
    assumptions.
\item \textbf{Chapter~\ref{ch:design}}: Presents the complete CBSR system architecture and
    formal design.
\item \textbf{Chapter~\ref{ch:implementation}}: Describes implementation details and
    performance optimizations.
\item \textbf{Chapter~\ref{ch:evaluation}}: Reports experimental results across five benchmarks.
\item \textbf{Chapter~\ref{ch:conclusion}}: Discusses limitations, summarizes contributions,
    and outlines future work.
\end{itemize}
